
We set out to measure whether improving AI quality makes users intellectually lazier, and found the opposite: returning users of the WildChat-1M platform compose progressively longer prompts over 2023--2024, not shorter ones. A Mann-Kendall trend test on 27,902 returning users (users appearing in $\geq 3$ distinct calendar months, April 2023--April 2024) yields $\tau$ = +0.744 (p = 0.0005, Hamed-Rao autocorrelation-corrected). This finding directly contradicts the central directional prediction of the Bidirectional Alignment Asymmetry (BAA) framework --- that as AI systems improve, human behavioral engagement should decline. Yet the result cannot be interpreted as simple refutation: the cohort design that makes behavioral signals detectable also systematically selects for high-engagement users, and this study contains no per-interaction measure of AI model quality, so the causal chain posited by BAA (AI improves $\rightarrow$ user behavior changes) is not directly tested. The result is a directional refutation in a specific cohort, not a causal one; the causal question remains empirically open.

The AI alignment literature has developed extensive frameworks for measuring how well AI systems align with human values and intent (Ouyang et al., 2022; Liang et al., 2022; Bai et al., 2022). Far less attention has been paid to the reverse question: how human behavior adapts in response to AI quality improvement over time. Shen et al. (2024) identify this bidirectional alignment measurement gap at the survey level, noting that while AI-to-human alignment is operationalized across dozens of benchmarks, human-to-AI behavioral alignment has no validated measurement framework. This asymmetry in the measurement infrastructure is the surface-level problem.

The deeper problem is mechanistic. If AI improvement reduces users' perceived need to carefully craft prompts, explicitly correct errors, or engage discriminatingly with AI outputs, then the behavioral feedback signals used to train and evaluate future models may become systematically less informative over time --- even as benchmark scores rise. This is the BAA conjecture: a divergence between AI benchmark improvement and human behavioral engagement, detectable computationally from existing interaction metadata without new annotation. We operationalize the BAA disengagement prediction as a negative monotonic trend in prompt token count ($\tau$ < 0). This operationalization follows from BAA's core mechanism (Shen et al., 2024 [UNVERIFIED]) but the specific proxy mapping --- that declining engagement manifests as shorter prompts --- is an extension of the framework and is not directly quoted from the source. Dell'Acqua et al. (2023) document qualitatively that AI-assisted consultants assign progressively simpler tasks to AI, suggesting reduced cognitive engagement, with effect sizes of approximately 0.3--0.5 SD. Perez et al. (2023) demonstrate that AI sycophancy suppresses explicit disagreement. Both lines of evidence suggest that declining behavioral engagement may be empirically real, but neither tests it at scale in public interaction logs.

The gap is twofold: no prior work has (1) empirically tested the BAA directional prediction at scale using public interaction data, or (2) validated a reusable measurement pipeline for detecting behavioral proxy trends in returning-user cohorts. This paper addresses both gaps and produces a negative result on the first that clarifies the second.

The key finding is that behavioral signals are computationally detectable in returning-user cohorts: prompt token count shows a strong, autocorrelated positive trend ($\tau$ = +0.744, ACF lag-1 = 0.634). This establishes that the measurement methodology works. But the direction of the trend refutes the BAA disengagement prediction. Whether the positive trend reflects user expertise gain, power-user self-selection, or a platform adoption effect is the central empirical question motivating future work.

We make four contributions:

\textbf{C1 (Empirical negative result).} To our knowledge, based on our survey of the literature, the first large-scale empirical test of the BAA disengagement directional prediction in public AI interaction logs. Returning WildChat-1M users do not show declining behavioral engagement: prompt length increases significantly ($\tau$ = +0.744, p = 0.0005), and no declining trend is detected for any proxy. All cited references are marked [UNVERIFIED] and require independent confirmation before this novelty claim can be asserted without qualification.

\textbf{C2 (Validated measurement pipeline).} A reusable analysis framework for behavioral trend detection in returning-user cohorts: WildChat-1M data loading, IP-hash cohort construction ($\geq 3$ monthly bins), Hamed-Rao Mann-Kendall testing, bootstrap confidence intervals (B = 1000, seed = 42), and figure generation.

\textbf{C3 (Documented infrastructure constraint).} The primary LMSYS Chatbot Arena dataset (\texttt{lmsys/chatbot\textbackslash \_arena\textbackslash \_conversations}) is access-gated on HuggingFace; the publicly available fallback dataset (\texttt{lmsys/lmsys-arena-human-preference-55k}) contains no timestamps. This is a binding constraint on any study requiring temporal vote entropy analysis over LMSYS preference data.

\textbf{C4 (Selection bias characterization).} We document and empirically quantify the returning-user cohort selection bias, and propose the specific comparison experiment needed to disentangle selection bias from genuine behavioral adaptation.

